% Part: first-order-logic
% Chapter: model-theory
% Section: partial-iso

\documentclass[../../../include/open-logic-section]{subfiles}

\begin{document}

\olfileid[fr]{mod}{bas}{pis}
\section{Isomorphismes partiels}

\begin{defn}
  Soient deux !!{structure}s $\Struct{M}$ et $\Struct{N}$.
  Un \emph{isomorphisme partiel} de $\Struct{M}$ dans $\Struct{N}$
  est une application partielle finie $p$ de $\Domain M$ dans
  $\Domain N$ qui vérifie, sur son domaine, les conditions
  d'isomorphisme de la \olref[iso]{defn:isomorphism} :
  \begin{enumerate}
  \item $p$ est !!{injective};
  \item pour chaque !!{constant}~$c$, si $p(\Assign{c}{M})$ est défini,
    alors $p(\Assign{c}{M}) = \Assign{c}{N}$;
  \item pour chaque !!{predicate} $n$-aire $P$, si $a_1$, \dots, $a_n$
    appartiennent au domaine de $p$, alors
    $\langle a_1, \dots, a_n\rangle \in \Assign P M$ si et seulement
    si $\langle p(a_1), \dots, p(a_n) \rangle \in \Assign P N$;
  \item pour chaque !!{function} $n$-aire $f$, si $a_1$, \dots, $a_n$
    et $\Assign f M (a_1, \dots,a_n)$ appartiennent au domaine de $p$,
    alors $p(\Assign f M (a_1, \dots,a_n))
    = \Assign f N (p(a_1), \dots, p(a_n))$.
  \end{enumerate}
  Dire que $p$ est finie signifie que $\dom{p}$ est fini.\footnote{Précision éditoriale : dans la condition sur une fonction, nous ajoutons que sa valeur dans $\Struct M$ doit appartenir au domaine de $p$, afin que l'expression $p(\Assign f M(\vec a))$ soit définie. Cette convention de correspondances finies n'exige pas que leur domaine soit une sous-structure close par les fonctions. La préservation de termes dont les valeurs sortent de ce domaine demande les prolongements d'une famille de va-et-vient.}
\end{defn}

Remarquons que l'application vide~$\emptyset$ est toujours un
isomorphisme partiel entre deux !!{structure}s quelconques.

\begin{defn}\ollabel{defn:partialisom}
  Deux !!{structure}s $\Struct{M}$ et $\Struct{N}$ sont
  \emph{partiellement isomorphes}, ce que l'on note
  $\Struct{M} \iso[p] \Struct{N}$, s'il existe un ensemble non vide
  $I$ d'isomorphismes partiels entre elles qui possède la propriété
  de \emph{va-et-vient} :
  \begin{enumerate}
  \item (\emph{Aller}) Pour tout $p \in I$ et tout
    $a \in \Domain M$, il existe $q \in I$ tel que $p \subseteq q$
    et que $a$ appartienne au domaine de $q$;
  \item (\emph{Retour}) Pour tout $p \in I$ et tout
    $b \in \Domain N$, il existe $q \in I$ tel que $p \subseteq q$
    et que $b$ appartienne à l'image de $q$.
  \end{enumerate}
\end{defn}

\begin{thm}\ollabel{thm:p-isom1}
  Si $\Struct{M} \iso[p] \Struct{N}$ et si $\Struct{M}$ et
  $\Struct{N}$ sont !!{enumerable}, alors $\Struct{M} \iso
  \Struct{N}$.
\end{thm}

\begin{proof}
  Puisque $\Struct{M}$ et $\Struct{N}$ sont !!{enumerable},
  écrivons $\Domain{M} = \{a_0, a_1, \ldots \}$ et
  $\Domain{N} = \{b_0, b_1, \ldots \}$. En partant d'un élément
  quelconque $p_0 \in I$, construisons une suite croissante
  d'isomorphismes partiels
  $p_0 \subseteq p_1 \subseteq p_2 \subseteq \cdots$ :
  \begin{enumerate}
  \item si $n+1$ est impair, donc $n = 2r$, la propriété d'aller
    fournit $p_{n+1} \in I$ tel que $p_n \subseteq p_{n+1}$
    et que $a_r$ appartienne au domaine de $p_{n+1}$;
  \item si $n+1$ est pair, donc $n+1 =2r$, la propriété de retour
    fournit $p_{n+1} \in I$ tel que $p_n \subseteq p_{n+1}$
    et que $b_{r-1}$ appartienne à l'image de $p_{n+1}$.
  \end{enumerate}
En posant
\[
p = \bigcup_{n\ge 0} p_n,
\]
on obtient un isomorphisme de $\Struct{M}$ sur $\Struct{N}$.\footnote{Correction éditoriale : le retour vise $b_{r-1}$, et non $b_r$ comme dans la source, qui omet ainsi $b_0$. Si l'un des domaines est fini, son énumération peut répéter des éléments ; les deux propriétés permettent alors de garder un prolongement déjà obtenu.}
\end{proof}

\begin{prob}
  Vérifier en détail que l'application $p$ construite dans le
  \olref[mod][bas][pis]{thm:p-isom1} est bien un isomorphisme.
\end{prob}

\begin{thm}\ollabel{thm:p-isom2}
  Supposons que $\Struct{M}$ et $\Struct{N}$ soient des !!{structure}s
  pour un !!{language} purement relationnel (ne contenant que des
  !!{predicate}s, sans !!{function} ni constante). Alors
  $\Struct{M} \iso[p] \Struct{N}$ entraîne
  $\Struct{M} \elemequiv \Struct{N}$.
\end{thm}

\begin{proof}
  Soit $I$ une famille de va-et-vient. On montre par induction sur
  les !!{formula}s que, si $p\in I$ envoie $a_i$ sur $b_i$
  pour $i =1$, \dots,~$n$, si les variables libres de $!A$
  sont parmi $x_1$, \dots, $x_n$, et si $s_1(x_i) =a_i$
  et $s_2(x_i) =b_i$, alors $\Sat{M}{!A}[s_1]$ si et seulement
  si $\Sat{N}{!A}[s_2]$. Les formules atomiques sont préservées
  par $p$; les connecteurs se traitent directement. Pour un
  quantificateur existentiel, on prolonge $p$ par la propriété
  d'aller afin d'y inclure un témoin dans $\Struct M$; la
  réciproque utilise la propriété de retour. Le cas $n=0$,
  appliqué aux énoncés, donne
  $\Struct{M} \elemequiv \Struct{N}$.\footnote{Précision éditoriale : l'induction porte sur tous les $p$ d'une même famille de va-et-vient $I$, et non sur une application partielle arbitraire comme le laisse entendre la source. La liste des variables libres est fixée ; les variables liées sont renommées au besoin.}
\end{proof}

\begin{rem}
En présence de !!{function}s, le résultat précédent reste vrai.
Pour $p\in I$, on le prolonge par l'aller afin d'inclure les valeurs
dans $\Struct M$ de toutes les constantes et de tous les sous-termes
d'une formule atomique donnée ; il suffit d'un nombre fini de prolongements.
Les conditions gardées assurent alors la correspondance des valeurs des termes et des formules atomiques. La même induction sur les formules s'applique.\footnote{Précision éditoriale : un $p$ isolé ne suffit pas ; cette justification utilise la famille de va-et-vient. Elle explicite le passage aux sous-structures engendrées évoqué par la source.}
\end{rem}

Le résultat précédent peut se décomposer en étapes : le nombre
de quantificateurs emboîtés dans !!a{formula} est lié au nombre
de prolongements d'isomorphismes partiels dont on a besoin.

\begin{defn}
  Pour toute !!{formula}~$!A$, son \emph{rang de quantification},
  noté $\QuantRank{!A} \in \Nat$, est défini récursivement
  comme le plus grand nombre de quantificateurs emboîtés dans
  $!A$. Deux !!{structure}s $\Struct{M}$ et $\Struct{N}$ sont
  \emph{$n$-équivalentes}, ce que l'on note
  $\Struct{M} \elemequiv[n] \Struct{N}$, si elles satisfont les mêmes
  énoncés de rang de quantification inférieur ou égal à~$n$.
\end{defn}

\begin{prop}\ollabel{prop:qr-finite}
  Soit $\Lang{L}$ un !!{language} relationnel fini, comportant un
  nombre fini de !!{predicate}s et de !!{constant}s, mais aucun
  !!{function} d'arité positive. Pour tout $n \in \Nat$ et tout
  ensemble fini fixé de variables libres, il n'existe, à
  équivalence logique près, qu'un nombre fini de !!{formula}s
  du premier ordre de $\Lang{L}$ dont les variables libres
  appartiennent à cet ensemble et dont le rang de quantification
  ne dépasse pas $n$. En particulier, cette conclusion vaut
  pour les !!{sentence}s.\footnote{Précision éditoriale : la source emploie ici « purement relationnel » tout en autorisant des constantes, et n'énonce la finitude que pour les énoncés. La preuve suivante a besoin de la version avec un ensemble fini fixé de variables libres. À chaque étape de l'induction, on ajoute une seule variable fraîche aux variables fixées, après renommage des variables liées.}
\end{prop}

\begin{proof}
  Par induction sur $n$ : pour un nombre fixé de variables,
  il n'y a qu'un nombre fini de formules atomiques distinctes
  dans ce langage; à chaque étape, les combinaisons booléennes
  et la quantification d'une variable ne produisent qu'un
  nombre fini de classes d'équivalence.
\end{proof}

\begin{defn}
  Étant donné !!a{structure}~$\Struct{M}$, soit
  $\Domain M^{<\omega}$ l'ensemble des suites finies d'éléments
  de $\Domain{M}$. Les lettres $\mathbf{a}, \mathbf{b},
  \mathbf{c}, \ldots$ désignent de telles suites. Si
  $\mathbf{a} \in \Domain{M}^{<\omega}$ et $a \in \Domain{M}$,
  alors $\mathbf{a}a$ désigne la \emph{concaténation} de
  $\mathbf{a}$ et de $a$.
\end{defn}

\begin{defn}
  Étant donné deux !!{structure}s $\Struct{M}$ et $\Struct{N}$,
  définissons par récurrence sur $n$ des relations
  $I_n \subseteq \Domain M^{<\omega} \times \Domain N^{<\omega}$
  entre suites de même longueur :
   \begin{enumerate}
   \item $I_0(\mathbf{a},\mathbf{b})$ si et seulement si
     $\mathbf{a}$ et $\mathbf{b}$ satisfont les mêmes
     !!{formula}s atomiques dans $\Struct{M}$ et $\Struct{N}$.
     Plus précisément, si ces suites sont de longueur $k$,
     si $s_1(x_i) = a_i$ et $s_2(x_i) = b_i$ pour
     $i=1$, \dots, $k$, et si les variables libres de
     la formule atomique $!A$ sont parmi $x_1$, \dots,~$x_k$,
     alors $\Sat{M}{!A}[s_1]$ si et seulement si
     $\Sat{N}{!A}[s_2]$.\footnote{Correction éditoriale : $k$ désigne la longueur des suites, indépendamment de l'indice de rang $n$. La source réemploie $n$ dans les deux rôles ; dans la définition de $I_0$, il faut bien considérer toutes les coordonnées de ces suites.}
   \item $I_{n+1} (\mathbf{a},\mathbf{b})$ si et seulement si,
     pour tout $a\in \Domain M$, il existe
     $b\in \Domain N$ tel que
     $I_n(\mathbf{a}a,\mathbf{b}b)$, et réciproquement.
   \end{enumerate}
\end{defn}

\begin{defn}
  Écrivons $\Struct{M} \approx_n \Struct{N}$ lorsque
  $I_n(\emptyseq,\emptyseq)$ est vrai pour $\Struct{M}$
  et $\Struct{N}$, où $\emptyseq$ est la suite vide.
\end{defn}

\begin{thm}\ollabel{thm:b-n-f}
  Soit $\Lang{L}$ un !!{language} purement relationnel.
  Si $I_n(\mathbf{a},\mathbf{b})$, alors pour toute
  formule $!A$ dont les variables libres figurent parmi
  les coordonnées de ces suites et telle que
  $\QuantRank{!A} \le n$, on a
  $\Sat{M}{!A}[\mathbf{a}]$ si et seulement si
  $\Sat{N}{!A}[\mathbf{b}]$; ici,
  $\mathbf{a}$ satisfait $!A$ si une assignation $s$
  telle que $s(x_i) = a_i$ satisfait $!A$.
  Si, de plus, $\Lang{L}$ est fini, la réciproque est vraie.
\end{thm}

\begin{proof}
  L'implication directe se montre par induction sur le rang
  et sur la construction de $!A$ : $I_n$ entraîne $I_0$
  en itérant les prolongements (les domaines des structures
  étant non vides); les formules atomiques relèvent donc de
  $I_0$. Les connecteurs préservent l'assertion et le cas
  existentiel utilise les deux clauses de va-et-vient de
  $I_{n+1}$.
  Pour la réciproque, raisonnons par induction sur $n$ en
  utilisant \olref{prop:qr-finite}, qui donne, pour un
  ensemble fixé de variables libres et une borne de rang,
  un nombre fini de classes de formules modulo
  l'équivalence logique.

  Pour $n=0$, si $\mathbf{a}$ et $\mathbf{b}$ satisfont
  les mêmes !!{formula}s sans quantificateur, elles
  satisfont les mêmes formules atomiques; donc
  $I_0(\mathbf{a},\mathbf{b})$.

  Supposons que $\mathbf{a}$ et $\mathbf{b}$ satisfassent
  les mêmes !!{formula}s de rang de quantification au plus
  $n+1$. Pour obtenir $I_{n+1}(\mathbf{a},\mathbf{b})$,
  il suffit, par hypothèse d'induction, de trouver,
  pour tout $a \in \Domain M$, un $b \in \Domain N$
  tel que $\mathbf{a}a$ et $\mathbf{b}b$ satisfassent
  les mêmes !!{formula}s de rang au plus $n$,
  puis de procéder dans l'autre sens.

  Pour $a \in \Domain M$ donné, posons $x=x_{k+1}$, avec $k$ la longueur de $\mathbf a$, et $\mathbf y=(x_1,\dots,x_k)$ ; les variables liées sont renommées au besoin. Soit $!T^a_n$
  l'ensemble des !!{formula}s $!B(x,\mathbf{y})$
  de rang au plus $n$ satisfaites par
  $\mathbf{a}a$ dans $\Struct{M}$. Cet ensemble peut
  être infini syntaxiquement, mais, à équivalence
  logique près, il possède un nombre fini de représentants.
  Leur conjonction constitue donc une seule !!{formula}\footnote{Correction éditoriale : la source dit que l'ensemble de toutes ces formules est fini. Il est généralement infini syntaxiquement ; seule la finitude des classes modulo l'équivalence logique autorise le choix d'une conjonction finie de représentants.}
  $!T^a_n(x,\mathbf{y})$. Comme, pour chaque classe,
  la formule ou sa négation est satisfaite par
  $\mathbf{a}a$, cette conjonction caractérise
  entièrement les formules de rang au plus $n$
  satisfaites par ce uplet. Dès lors,
  $\mathbf{a}$ satisfait
  $\lexists[x][!T^a_n(x,\mathbf{y})]$,
  dont le rang ne dépasse pas $n+1$.
  Par hypothèse, $\mathbf{b}$ satisfait la même
  !!{formula} dans $\Struct{N}$ : il existe donc
  $b \in \Domain N$ tel que $\mathbf{b}b$ satisfasse
  $!T^a_n$. Les uplets $\mathbf{a}a$ et
  $\mathbf{b}b$ satisfont alors exactement les mêmes
  !!{formula}s de rang au plus $n$.
  Le même raisonnement, en échangeant $\Struct{M}$
  et $\Struct{N}$, donne pour tout $b \in \Domain N$
  un $a\in \Domain M$ ayant cette propriété.
  La démonstration est achevée.
\end{proof}

\begin{cor}\ollabel{cor:b-n-f}
  Si $\Struct{M}$ et $\Struct{N}$ sont des !!{structure}s
  purement relationnelles dans un !!{language} fini, alors
  $\Struct{M} \approx_n\Struct{N}$ si et seulement si
  $\Struct{M} \elemequiv[n] \Struct{N}$. En particulier,
  $\Struct{M} \elemequiv \Struct{N}$ si et seulement si,
  pour tout $n$, $\Struct{M} \approx_n \Struct{N}$.
\end{cor}

\end{document}
